\backhead{FURTHER EXERCISES}{khred}
\begin{exlist}
\item How many runs does a two-level full factorial of eight factors need per soil? A Latin hypercube of 50 runs covers each factor in how many strata?\exans{256; 50}
\item A bootstrap of 4999 resamples gives the 2.5 and 97.5 percentiles of a constant. Which resamples (by rank) are these?\exans{the 125th and the 4875th}
\item In leave-one-soil-out validation of ten soils, how many fits are made, and on how many soils is each fitted?\exans{10; 9}
\item Two laws have LOSO RMSE 0.058 and 0.111. By what percentage is the first better?\exans{48 per cent}
\item A constant has the interval 58 to 105 around 71.5. Is the interval symmetric? Find the distances below and above.\exans{no; 13.5 and 33.5}
\end{exlist}
\backhead{ESSAY QUESTIONS}{khred}
{\small\itshape Answer each question in about one page, with equations, a sketch and at least one number from this chapter.}\par\vspace{3pt}
\begin{exlist}
\item Describe the design of the simulation experiments of the study and the reasons for each design.
\item Discuss leave-one-out validation and why it is the main test of the laws of this book.
\item Explain the bootstrap and what its intervals mean.
\item Discuss the verification of SWAFL against textbook examples and held-out runs.
\item Explain how the studies of the project fit together into one argument.
\item Choose one study and write a critical review of its methods and conclusions.
\end{exlist}
\dykband{EXERCISES IN SWAFL}
{\small Open SWAFL (desktop or online) with the benchmark project. The chapter of the \href{https://github.com/halilunaziru73-creator/Coupled-Water-Fluxes-Pore-Scale-Dynamics-Precision-Irrigation-Mulching-Feedback-Loops/raw/main/Software/v1.1.0/SWAFL_Desktop/docs/manual_v1.1.0/SWAFL_Manual_v1.1.0.pdf}{SWAFL manual} that explains each task is given in brackets; record your values in a table and compare them with what is expected (\emph{Check}).}\par\vspace{3pt}
\begin{exlist}
\item {\color{khblue}[Manual, Chapter 28]} Design a Daisy experiment with a Latin hypercube of 20 runs and inspect the plan.\par\noindent{\small\itshape Check: each factor is sampled in 20 strata.}
\item {\color{khblue}[Manual, Chapter 29]} Open the best combinations and choose one for your field.\par\noindent{\small\itshape Check: the plan of the combination is shown.}
\item {\color{khblue}[Manual, Chapter 30]} Read the interactions of an experiment and identify the strongest.\par\noindent{\small\itshape Check: usually between water and nitrogen.}
\item {\color{khblue}[Manual, Chapter 18]} Use the search for the most efficient inputs and compare with the best combination of Daisy.\par\noindent{\small\itshape Check: the two agree in direction.}
\end{exlist}
\backhead{OBJECTIVE TYPE QUESTIONS}{khred}
\begin{exlist}
\item A two-level full factorial of six factors has\opt{12}{36}{64}{720} runs
\item In leave-one-soil-out validation of ten soils each fit uses\opt{1}{5}{9}{10} soils
\item A bootstrap draws soils\opt{without replacement}{with replacement}{in order}{once each}
\item A Latin hypercube samples each factor\opt{at two levels}{once in each stratum}{at random only}{at its mean}
\item An overfitted law has\opt{a low LOSO error}{a low in-sample error and a high LOSO error}{no parameters}{a high in-sample error}
\end{exlist}
\objkey{1.\ (c),\quad 2.\ (c),\quad 3.\ (b),\quad 4.\ (b),\quad 5.\ (b)}
